\subsection{Step 8: extend the chain while retaining numerical precision}
\label{sec:extended-production}
\begin{equation}\mathcal D=\{2\operatorname{round}[30\,100^{j/48}]:j=0,\ldots,48\},\qquad \mathcal R=\{200+4j:j=0,\ldots,200\}.\label{eq:dense-grid}\end{equation}
\begin{align}
\mathcal N_{\rm ising}&=\mathcal N\cup\mathcal D\cup\mathcal R\cup\{4096+256j:j=0,\ldots,16\}\nonumber\\
&\quad\cup\{152,1024,1118,1146,1232,1356,1492,1528,1604,1644\}\nonumber\\
&\quad\cup\{1808,1852,1972,2048,2190,2354,2412,2654,2890,3000\}\nonumber\\
&\quad\cup\{3140,3216,3294,3436,3500,3552,3606,3660,3784,3854\}\nonumber\\
&\quad\cup\{3926,4000,4158,4222,4286,4426,4670,4734,4798,5000\}\nonumber\\
&\quad\cup\{5246,5310,5542,5758,6270,6526,6782,7038,7294,7550\}\nonumber\\
&\quad\cup\{7806,8000,8016,8032,8048,8064,8080,8096,8112,8128\}\nonumber\\
&\quad\cup\{8144\},\nonumber\\
\mathcal N_{\rm boson}&=\mathcal N\cup\mathcal D\cup\{100,104,168,292,322,354,428,472,512,762\}\nonumber\\
&\quad\cup\{924,1024,1118,1356,1492,1644,1808,2048,2412,2922\}\nonumber\\
&\quad\cup\{3540,3750,3916,4096,4234,4376,4608,4720,4834,5000\}\nonumber\\
&\quad\cup\{5222,5626,5810,6250,6466,6688,6918,7156,7402,7656\}\nonumber\\
&\quad\cup\{7920,8192,8436,8688,8948,9216,9462,9714,9974,10240\}\nonumber\\
&\quad\cup\{10488,10740,10998,11264,11764,12288,12800,13312,13814,14336\}\nonumber\\
&\quad\cup\{14840,15360,15864,16384\}.
\label{eq:extended-grid}\end{align}
All indices $j$ in these set definitions are integers within the stated bounds.

The earlier grid in Eq.~\eqref{eq:large-grid} remains a reference subset.
The active mixture fits use Eq.~\eqref{eq:extended-grid}; no intervals
longer than ten sites are added. The primary-to-primary sanity check
in Appendix~\ref{sec:primary-relative-entropy} retains its original
$N\leq256$ grid. The methods below change how the same finite-chain
states and entropies are evaluated, not their definitions.

\subsubsection{Fourier construction and local Ising transition solves}
At criticality the Majorana matrix is uniform apart from its boundary
twist, so its eigenmodes are known Fourier waves. Let $L=2N$ count
Majoranas, with indices $j=0,\ldots,L-1$, and set
$k_m=2\pi(m+t)/L$. The twist is $t=1/2$ for $I,\eps$ and $t=0$
for $\sigma$. Define $w_m=-\operatorname{sgn}(\sin k_m)$.
For the energy state reverse $w_m$ at $m=0,N-1,N,L-1$.
For the spin state set the two zero-mode weights $w_0,w_N$ to zero
and supply their parity explicitly. Then
\begin{equation}
 (\Gamma_a)_{ij}=\frac{i}{L}\sum_{m=0}^{L-1}w_m
 e^{ik_m(i-j)}+\begin{cases}
 [(-1)^i-(-1)^j]/L,&a=\sigma,\\0,&a=I,\eps.
 \end{cases}
 \label{eq:fourier-covariance}
\end{equation}
The imaginary notation in the Fourier sum yields a real antisymmetric
matrix. The last term is $vu^T-uv^T$, with $u_j=L^{-1/2}$ and
$v_j=(-1)^jL^{-1/2}$; it selects the same odd Ramond state as
Sec.~\ref{sec:large-ising}. This is an exact finite-chain construction,
using the free-fermion mode description of Refs.~\cite{Montes2017,PeschelEisler}.
Its energies are
\begin{equation}
 E_I=-2\csc\frac{\pi}{2N},\quad
 E_\sigma=-2\cot\frac{\pi}{2N},\quad
 E_\eps=E_I+8\sin\frac{\pi}{2N}.
 \label{eq:fourier-energies}
\end{equation}
A fast Fourier transform applies $\Gamma_a$ to a vector in
$O(N\log N)$ work. Thus the Hamiltonian eigensolve is unnecessary.

For the anchored transition in Sec.~\ref{sec:complement-ising}, let
$\Gamma'_b$ be the covariance after applying the local anchor.
Form $K=\id-\Gamma_a\Gamma'_b$ by applying the Fourier transform to
batches of columns. Factor the real matrix $K$ once with pivoted LU
(lower and upper triangular factors). Let $E_A$ be the first $2\ell$
columns of the identity. Solve only
\begin{equation}
 KZ=(\id-i\Gamma_a)E_A,\qquad
 \Gamma_{ab,A}=i\{E_A^T(\id-i\Gamma'_b)Z-\id_{2\ell}\}.
 \label{eq:local-transition-solve}
\end{equation}
The real and imaginary right-hand sides reuse the same LU factors.
The anchor overlap is $|\det K|^{1/4}/2^{N/2}$; its logarithm is
obtained from the LU diagonal. Reconstruct the local Gaussian operator
and restore the anchor, exactly as in the earlier Gram calculation.
We never construct a full complex transition covariance or its inverse.
The remaining dense factorization costs $O(N^3)$ arithmetic and the
stored real matrices cost $O(N^2)$ memory. The interval density matrices
still have dimension at most $2^{10}=1024$.

\subsubsection{Direct evaluation of small entropy remainders}
\label{sec:stable-entropy}
At very small fractions, subtracting $\log2-D^B$ or two almost equal
entropy traces loses relative precision. Instead, evaluate their
remainder directly with matrix-function divided differences
\cite{Higham}. The following identity is an algebraic numerical device,
not a CFT expansion or a new analytical entropy coefficient.

Set $f(t)=t\log t$, with $f(0)=0$. In a basis where a positive reference
matrix is $D=\operatorname{diag}(d_i)$, write
$H=D+E=V\operatorname{diag}(\lambda_k)V^\dagger$.
The first-order matrix derivative has diagonal $f'(d_i)E_{ii}$.
For the remainder diagonal, use
\begin{equation}
 \mathcal R_i=[f(H)-f(D)]_{ii}-f'(d_i)E_{ii}
 =\sum_k |(V^\dagger E)_{ki}|^2\,
 \frac{\lambda_k\log(\lambda_k/d_i)-\lambda_k+d_i}
      {(\lambda_k-d_i)^2}.
 \label{eq:stable-remainder}
\end{equation}
To verify it, expand $[f(H)]_{ii}=\sum_k|V_{ik}|^2f(\lambda_k)$,
use $\sum_k|V_{ik}|^2=1$ and
$E_{ii}=\sum_k|V_{ik}|^2(\lambda_k-d_i)$, and factor the scalar
Taylor remainder. Finally,
$(\lambda_k-d_i)V_{ik}=(EV)_{ik}$ converts its squared factor to
the numerator weight in Eq.~\eqref{eq:stable-remainder}.
At equal arguments the quotient is $1/(2d_i)$.
With $u=(\lambda-d)/d$, compute it as
\begin{equation}
 \frac1d\frac{(1+u)\log(1+u)-u}{u^2}
 =\frac1d\sum_{n=0}^{\infty}\frac{(-u)^n}{(n+1)(n+2)}.
 \label{eq:stable-scalar}
\end{equation}
The code uses ten terms only for $|u|<0.02$, where the omitted
tail is below double-precision accuracy, and the closed expression
otherwise. This series stabilizes a scalar matrix-function evaluation;
it is unrelated to a small-interval power-law fit.

For $D$ the eigenvalue matrix of the mixture on $A$ and $H$ the source
state in that basis, equal unit traces give
$S(\rho\|\mu)=\sum_i\mathcal R_i$. For the complement, first define
the local transition $T_{ab}^A=\Tr_B\ket a\bra b$.
The unweighted Gram matrix has blocks
$H=\left(\begin{smallmatrix}\rho_a^A&T_{ab}^A\\
(T_{ab}^A)^\dagger&\rho_b^A\end{smallmatrix}\right)$.
Its relation $H=2\mathcal G^*$ to the weighted Gram matrix is derived
in Sec.~\ref{sec:complement-gram}. Transform the two diagonal blocks
to their separate Schmidt eigenbases and let $D$ contain their
eigenvalues on its diagonal.
Then $E$ contains only the off-diagonal transition blocks, so its
first-order diagonal vanishes. If $P_a,P_b$ select the two state blocks,
the Gram entropy identity gives the deficit
$\delta_{a\mid b}\equiv\log2-D^B_{a\mid b}$ directly:
\begin{equation}
 \delta_{a\mid b}=\Tr P_a[f(H)-f(D)]
       =\sum_{i\in a}\mathcal R_i,\qquad
 \delta_{b\mid a}=\sum_{i\in b}\mathcal R_i.
 \label{eq:stable-deficit}
\end{equation}
Complex conjugating the entire Gram matrix preserves these real block
traces. Thus even a deficit much smaller than machine precision
relative to $\log2$ can be stored accurately on its own.

The identities are exact on a positive support. In double precision
we retain reference eigenvalues above $\tau=10^{-15}$ in each
symmetry sector. The discarded tails are a numerical approximation;
they are checked by varying $\tau$ to $10^{-14}$ and $10^{-16}$ at
every new size for $\ell=4,10$. We additionally compare with the original
trace-log and Gram formulas and with a high-precision, noncommuting
two-dimensional example. These checks, together with the larger-size
state-contraction tests, are reported in Appendix~\ref{sec:scaling-checks}.

\subsubsection{Boson scaling and the practical boundary}
The boson continues to use the exact sparse moment and small Pfaffian
reductions in Sec.~\ref{sec:boson-contraction} and
Appendix~\ref{sec:complement-pf-details}. The dense high-precision
minors depend on $\ell$, not on $N$. At the added sizes the moment
algebra uses 120 decimal digits before its final conversion to double
precision. Independent sampled entries use 160 digits and direct
pivoted Pfaffians instead of cached recursive minors.
Although the physical moment matrices are sparse, the current Python
pivot searches scan remaining indices. Their bookkeeping can cost
$O(N^2)$ time, in addition to the arithmetic and the $O(N\ell)$ stored
right-hand sides. This becomes the main practical obstacle before the
$2^\ell$ interval dimension changes. The measured limits and timings
below refer to this implementation on this machine; they are not
mathematical maximum chain lengths.
\begin{longtable}{lrrrr}
\caption{Measured largest completed chains, with three diagonal states, three transitions, both directions, and every $\ell=4,\ldots,10$. Wall time includes both interval calculations and endpoint cutoff checks. Memory is the process high-water resident set (GiB), not the size of one matrix; it may include allocations retained from preceding stages. The host has 32 GiB RAM. These are validated endpoints, not absolute mathematical limits.}
\label{tab:scaling-resources}\\\toprule
CFT & Largest $N$ & Seconds & Peak GiB & Minimum retained fraction\\\midrule
\endfirsthead
\toprule
CFT & Largest $N$ & Seconds & Peak GiB & Minimum retained fraction\\\midrule
\endhead
\bottomrule\endfoot
ising & 8192 & 93.46 & 9.0882 & $7.324\!\times\!10^{-4}$\\
boson & 16384 & 506.34 & 0.73204 & $3.662\!\times\!10^{-4}$\\
\end{longtable}
For Ising, doubling the largest tested $N$ would multiply the dominant stored matrix sizes by four, giving an estimated 36.4 GiB with the present allocation pattern, above this host\textquotesingle s RAM. The dense LU arithmetic would grow by about eight. For the boson, sparse index scans still permit further growth but become costly: doubling $N$ can quadruple that bookkeeping. Such extrapolations are estimates, not executed benchmarks. No larger size is claimed as validated.

